\documentclass[10pt]{article}

\usepackage[utf8]{inputenc}

\usepackage[T1]{fontenc}

\usepackage{amsmath}

\usepackage{amsfonts}

\usepackage{amssymb}

\usepackage[version=4]{mhchem}

\usepackage{stmaryrd}

\usepackage{graphicx}

\usepackage[export]{adjustbox}

\graphicspath{ {./images/} }

\usepackage{mathrsfs}



\begin{document}

конечного числа законов сохранения означает, что при рассеянии сохраняются количества частиц каждого типа; $n$-частичная $S$-матрица сводится к парным $S$-матрицам. С помощью интеграла по траекториям можно вычислить квантовые поправки к массам и к квазиклассическоӥ $S$-матрице солитонов. Одним из нетривиальных свойств указанной модели является появление целого спектра частиц (солитонов), в то время как лагранжиан теории содержит только одно поле. Кіроме того, в приближении слабого взаимодействия (т е когда $\gamma$ мало) солитоны - тяжелые частицы и сильно взаимодействуют.



Jum.: [1] A b l o w itz M. [и дp.], "Phys. Rev. Lett.",



\includegraphics[max width=\textwidth, center]{2023_07_08_1b0cc69c594d48a31546g-1(2)}

174; [3] и X ж e, «Tp. Матем. ин-та АH СCCP», 1976, т. 142, с. $254-66 ;[4]$ К о з е л В. А., К о т л я р о в В. П., "Докл АН УCCP, сер. А.», 1976, ㅊ 10, с. 878-81; [5] К о р е-

п и н В. Е., Ф а д д е е в Л.' Д., «Теоретич. и матем. Физика», 1975 , т. 25, № 2, c. $147-63$; [6]'B i a n c h i L., Lezioni di geometria differenziale, v. 2, pt 1-2, Pisa, 1923-24; [7] $\Phi_{\text {и н и- }}$ Кㅇ в С. П., Изгибание на главном основании..., М.-- Ј., 1937; т. 19, №5-6, с. 883-901.



СИНУСОВ ТЕОРЕМА: для произвольного треугольника со сторонами $a, b, c$ и противолежащими им углами $A, B, C$ имеют место соотношения



\[

\frac{a}{\sin A}=\frac{b}{\sin B}=\frac{c}{\sin C}=2 R

\]



где $R$ - радиус описанного круга. Ю. А. Горъков.



СИНУСОИДА - график функции $y=\sin x$ (см. рис.). С.- непрерывная кривая с периодом $T=2 \pi$. Пересе-



\begin{center}

\includegraphics[max width=\textwidth]{2023_07_08_1b0cc69c594d48a31546g-1}

\end{center}



чения с осью $O x$ - точки $(k \pi, 0)$; они же - точки перегиба с углом $\pm \pi / 4$ наклона к оси $O x$; экстремумы $\left((k+1 / 2) \pi,(-1)^{k}\right)$.



График функции $y=\cos x=\sin (x+\pi / 2)-\kappa$ о с и н ус о и д а - представляет собой С., сдвинутую влево на $\pi / 2$. Пересечения косинусоиды с осью $O x$ - точки $((k+1 \%) \pi, 0) ;$ экстремумы $\left(k \pi,(-1)^{k}\right)$.



Многие колебательные процессы описываются периодич. функцией вида $y=a \sin (b x+c)$, где $a, b$ и $c-$ постоянные. График этой функции (т. н. о б ша а я с и н у с о и д а, по сравнению с графиком функции $y=\sin x \quad($ б б ы к о о в е н а я с и н у с о и д а) вытянут вдоль оси $O y$ в $|a|$ раз, сжат вдоль оси $O x$ в $b$ раз, сдвинут влево на отрезок $c / b$ и при $a<0$ зеркально отражен относительно оси $O x$; период $T=2 \pi / b$, пересечения с осью $O x-$ точки $\left(\frac{k \pi-c}{b}, 0\right)$; әкстремумы $\left(\frac{(k+1 / 2) \pi-c}{b},(-1)^{k} a\right)$. ю. А. Горъков.



СИНУСОИДАЛЬНАЯ СПИРАЛЬ - плоская кривая, уравнение к-рой в полярных координатах имеет вид



\[

\rho^{m}=a^{m} \cos m \varphi \text {. }

\]



При рациональном $m$ С. с. является алгебраич. кривой. В частности, при $m=1$ - окружность, при $m=-1-$ равносторонняя гипербола, при $m=1 / 2$ - кардиоида, при $m=-1 / 2$-парабола.



В общем случае для $m>0$ С. с. проходит через полюс и полностью помещается внутри круга радиуса $a$. При отрицательном $m$ радиус-вектор кривой может принимать сколь угодно большие значения и не праходит через полюс. С. с. симметрична относительно полярной оси, при рациональном $m=p / q$ (где $p$ и $q$ взаимно простые числа) имеет $p$ осей симметрии, проходящих через полюс. При целом положительном $m$ радиус-вектор С. с. является периодич. функцией с периодом $2 \pi / \mathrm{m}$. При изменении $\varphi$ от 0 до $2 \pi$ кривая состоит из $m$ лепестков, каждый из которых располагается внутри угла, равного $\pi / m$. Полюс в этом случаекратная точка (см. рис.).



\includegraphics[max width=\textwidth, center]{2023_07_08_1b0cc69c594d48a31546g-1(1)}

$\Pi p и$ дробном положительном

$m=p / q$ кривая состоит из $p$ пересекающихся лепестков. При целом отрицательном $m$ C. c. состоит из $|m|$ бесконечных ветвей, к-рые могут быть получены инверсией спирали с $m^{\prime}=-m$.



Лит [1] С а в е ло в А. А., Плоские кривые, М., 1960. СИНУС-ПРЕОБРАЗОВАНИЕ ФУРБЕ - см. Фурье преобразование.



$A$-СИСТЕМА - счетно ветвящаяся система множеств, т. е. семейство $\left\{A_{n_{1} \ldots n_{k}}\right\}$ подмножеств множества $X$, занумерованных всеми конечными последовательностями натуральных чисел. $A-\mathrm{C} .\left\{A_{n_{1} \ldots n_{k}}\right\}$ наз. p eг у л я р но й, если $A_{n_{1} \ldots n_{k}} n_{k+1} \subset A_{n_{1} \ldots n_{k}}$. Последовательность $A_{n_{1}}, \ldots, A_{n_{1} \ldots n_{k}}, \ldots$ элементов $A$-С., занумерованных всеми отрезками одной и той же бесконечной последовательности натуральных чисел, наз. ц е п ь ю этой $A$-С. Пересечение всех элементов цепи наз. ее я д р о м, а объединение ядер всех цепей $A$-С.ядром этой $A$-С., или р е з у ли т а т о м $A$-о п е р ац и и, примененной к этой $A-C .$, или $A-\mathrm{м}$ в о ж е с тв о м, порожденным этой $A$-С. Всякая $A$-С. может быть регуляризована без изменения ядра (достаточно положить $\left.A_{n_{1}}^{\prime} \ldots n_{k}=A_{n_{1}} \cap \ldots \cap A_{n_{1}} \ldots n_{k}\right)$. Если $\mathscr{M}-$ нек-рая система множеств, то ядра $A$-С., составленных из әлементов системы $\mathscr{M}$, наз. $A-$ м н о ж е с т в а м и, порожденными системой $\mathscr{M} . A$-множества, порожденные замкнутыми множествами топологич. пространства, наз. $A$-множествами әтого пространства.



Лит.: [1] А л е к с а н д р о в П. С., Введение в теорию множеств и общую топологию, М., 1977; [2] К у р а т о в с к и й К-СИСТЕМА $\left\{T^{t}\right\}$ - такой измеримый поток (Kпо т о к) или каскад (K-к а с к а д) в Лебега пространстве, что существует измеримое разбиение $\xi$ фазового пространства со следующими свойствами: а) оно в о зр а с т а ю щ е е (в более старой терминологии - и нв а р и а н т о о е) относительно $\left\{T^{t}\right\}$, то есть $T^{t \xi}$ мельче $\xi \bmod 0$ при $t>0 ;$ б) оно является д в у с т оро нни и о б р а з у ю ш и м для $\left\{T^{t}\right\}$, т. е. единственным mod 0 измеримым разбиением, к-рое мельче $\bmod 0$ всех $T^{t \xi}$, является разбиение на точки; в) единственным mod 0 измеримым разбиением, к-рое крупнее $\bmod 0 \operatorname{scex} T^{t} \xi$, является тривиальное разбиение, единственный элемент к-рого - все фазовое пространство.



Автоморфизм пространства с мерой, итерации к-рого образуют $K$-каскад, наз. $K-2$ в то м о р ф и з мо м. Если $\left\{T^{t}\right\}-K$-С., то все $T^{t}$ с $t \neq 0$ являются $K$-автоморфизмами; обратно, если для измеримого потока или каскада $\left\{T^{t}\right\}$ хоть одно $T^{t}$ является $K$-автоморфизмом, то $\left\{T^{t}\right\}-K-C$. $K$-система обладает сильными әргодич. свойствами: положительная әнтрошия, әргодичность, перемешивание всех степеней, счетнократ-





\end{document}